\section{Conclusion}

This paper presented a vehicle-cloud collaborative predictive planning and anti-rollover control framework for autonomous tank trucks. The framework follows a nested safety logic: the cloud-side outer loop anticipates interaction-induced risks through Graphflow-based prediction and reinforcement-learning decision-making, the vehicle-side planning loop provides contingency fallback when the cloud reference becomes locally unsafe, and the high-frequency control loop suppresses roll and liquid-sloshing state risks through multi-constraint MPC. In this structure, the cloud trajectory is predictive guidance rather than a direct command, and the vehicle remains responsible for real-time feasibility and anti-rollover safety.

Validation was conducted through traffic simulation, vehicle-liquid co-simulation, and scaled tank-truck experiments. In the G15 highway traffic simulation, Graphflow provided long-horizon risk-field prediction and the cloud-side policy improved safety-related scores over a single-vehicle baseline while maintaining comparable efficiency. The vehicle-side planner generated feasible contingency behaviors in representative local conflict cases. In co-simulation, the proposed steering-braking MPC kept the equivalent LTR within the soft safety boundary in high-risk left-turn, single-lane-change, and double-lane-change scenarios where pure tracking or late differential braking failed or approached rollover. In scaled experiments, the controller reduced the maximum roll angle by 42\% and achieved real-time execution with an average solution time of 15.4 ms.

Several limitations remain. Full-scale real-vehicle validation has not yet been conducted, and the influence of communication delay, packet loss, and cloud-edge deployment constraints should be studied in closed-loop experiments. The current traffic validation focuses on highway and interchange scenarios; more complex mixed-traffic, ramp, construction-zone, and adverse-weather cases should be added. Future work will also improve quantitative ablation coverage and investigate deployment-efficient cloud-edge implementations for real hazardous-liquid transportation systems.
